\section{总结}

本文演示了一种从格点计算抽取部分子分布的新方法。
它建立在文献 \cite{Radyushkin:2017cyf} 所阐述的思想之上。

首先，我们把一般性等时矩阵元处理为 Ioffe 时间 $\nu = Pz_3$
与距离 $z_3$ 的函数 ${\cal M} (\nu, z_3^2)$。
下一步的想法是构造 Ioffe 时间分布 ${\cal M} (\nu, z_3^2)$
与静止系密度 ${\cal M} (0, z_3^2)$ 之比
${\mathfrak  M} (\nu, z_3^2) \equiv {\cal M} (\nu, z_3^2)/{\cal M} (0, z_3^2)$。


我们的格点计算清楚地表明，静止系函数的 $z_3$ 依赖中存在线性成分，
它可以归因于规范链接所生成的预期行为 $Z(z_3^2) \sim e^{-c|z_3|/a}$。
接着我们观察到，对计算中所用的全部 6 个动量 $P$ 取值，
比值 ${\cal M} (Pz_3 , z_3^2)/{\cal M} (0, z_3^2)$ 对 $z_3$ 均呈高斯型行为。
这意味着进入 ${\mathfrak  M} (Pz_3, z_3^2)$ 比值分子与分母的
$Z(z_3^2)$ 因子已如预期被消去。

不过，没有任何 {\it 先验}原理保证比值 ${\cal M} (\nu, z_3^2)/{\cal M} (0, z_3^2)$
分子与分母之间的非对数型残余 $z_3^2$ 依赖一定相互抵消。
如有需要，这类 $z_3^2$ 依赖可以通过系统的拟合程序予以去除，
从而在 $z_3^2=0$ 极限下抽取 Ioffe 时间 PDF。

然而我们发现，当把 ${\mathfrak  M} (\nu, z_3^2)$ 实部与虚部的数据
作为 $\nu$ 与 $z_3$ 的函数作图时，
它们都非常接近各自的普适函数。
这一观察表明 ${\cal M} (\nu, z_3^2)$ 的软 $z_3^2$ 依赖部分
已被静止系密度 ${\cal M} (0, z_3^2)$ 抵消。
该现象对应 TMD ${\cal F} (x, k_\perp^2)$ 中 $x$ 依赖与 $k_\perp$ 依赖的因子化。

虽然这一支持因子化性质的证据本身就是一个重要结果，
我们要强调的是，我们的方法并{\it 不依赖}于因子化。
它基于比值 ${\cal M} (\nu, z_3^2)/{\cal M} (0, z_3^2)$ 的使用；
其残余的软 $z_3^2$ 依赖可以被系统地分析与拟合，
因而 $z_3^2$ 极限能够以可控的方式取得。

幸运的是，在我们目前的统计与系统误差范围内，
数据并未表现出可见的多项式 $z_3^2$ 依赖。
在今后的工作中，我们打算仔细研究残余的多项式 $z_3^2$ 效应，
并把它们纳入这里引入的格点 PDF 抽取方法之中。

此外，我们还验证了：对于小的 $z_3 \leq 4a$，
残余 $z_3$ 依赖可以用 $\alpha_s/\pi =0.1$ 对应的微扰演化解释。
我们把这些小 $z_3$ 数据点演化到 $z_3=2a$ 标度（对应 $\mu^2=1$ GeV$^2$），
演化后的数据更好地逼近了 ${\cal M}$ 实部与虚部的普适曲线，
支持观测到微扰演化的论断。

因此，这些普适曲线的 $ \nu \lesssim 4$ 部分
可视为对应 $\mu=1$ GeV 标度。
其余数据点对应 $z_3 >4a$，
形式上应视为对应 $\mu \lesssim 0.3$ GeV 的标度。
所有这些数据点基本上都落在同一条普适曲线上，
表明演化在这样的标度处已经停止。
我们把这条“低归一化点”曲线与三个演化到 $\mu =1$ GeV 标度的全局拟合进行了比较，
观察到我们关于价 $u_v(x)- d_v (x)$ 分布的曲线 (\ref{qV})
在 $x \to 1$ 处呈现 $(1-x)^3$ 行为，符合通常预期；
而且在相当小的 $x$ 范围内都比较贴近 NNPDF31、
尤其是 MSTW 的 NNLO 全局拟合。


不过，我们的曲线在 $x <0.1$（NNPDF31 情形）
与 $x <0.05$（MSTW 情形）处与全局拟合强烈偏离。
但 PDF 的形状会受到微扰演化的影响。
为展示这些效应的范围，我们把所有 $z_3 \leq 10a$ 的数据点
演化到对应 $\mu =1 $ GeV 的普适标度 $z_0=2a$，
再把所得 PDF 进一步演化到 $\mu =2 $ GeV——
即全局拟合的标准参考标度。
我们的最终曲线与这些拟合相当接近，
这表明在格点结果与实验数据的比较中，
微扰演化扮演着重要角色。
同样地，要在足够宽的 Ioffe 时间参数 $\nu$ 区间内为微扰演化方程的使用提供依据，
还需要更小的格距。


数据还表明存在非零的{\it 正}反夸克分布 $\bar q (x) = \bar u(x) - \bar d(x)$。
它使 $q(x)$ 的 $x$ 积分改变 7\%，且具有 $\sim x (1-x)^3$ 行为。
由于采用的是淬火近似，这些反夸克来自“连接图”（connected diagrams）。
因此可以预期比值 $\bar u/\bar d$ {\it 必须}跟随质子的味组分，
即 $\bar u/\bar d \sim 2$、$\bar u >\bar d$。
我们的数据与该预期一致。

本研究具有探索性质，其主要目标是发展基于文献 \cite{Radyushkin:2017cyf} 思想的
格点 PDF 抽取技术。
我们的结果表明，所提出的基本方法在从格点 QCD 获得可靠 PDF 方面潜力巨大。
今后我们将改进方法，以纳入演化并在 Ioffe 时间分布抽取中控制残余的多项式 $z_3^2$ 效应。

为此，显然需要更小的格距以及更大的核子动量范围。
此外，我们还需研究有限体积效应，
并纳入 π 介子质量更接近物理点的动力学费米子。
我们计划在未来的工作中处理所有这些问题。
